\section{Removing amplification under nuclearity}\label{sec:nuclear}

We now remove the fixed amplification under nuclearity.  This proof
uses the graph construction of Section~\ref{sec:graph}: a finite tuple
of matrix entries rules out the graph's failure relations, and the
resulting shrinking operation allows us to join rich positive
elements detected by different representations.  Its second output
need not be rich; the merger retains \(n\) orthogonal copies so that
\(\propF\) makes their sum rich.  A finite ideal argument then reaches every positive
element.  The proof of the following theorem is
independent of the exact localization argument in
Sections~\ref{sec:exactness}--\ref{sec:localization}.

\begin{theorem}[Removing amplification under nuclearity]
\label{thm:nuclear-to-individual}
Let \(C\) be a nuclear complex \(C^*\)-algebra.  Suppose that one
integer \(n\ge1\), fixed for the whole algebra, makes
\(h^{\oplus n}\) properly infinite for every \(h\in C_+\).
Then every positive element of \(C\) is properly infinite.
\end{theorem}

Theorem~\ref{thm:nuclear-to-individual} assumes neither separability
nor a unit.  We keep the fixed integer \(n\) throughout this section.
In the tuple construction it controls the number of graph entries,
while the dimensions of the finite matrix approximations may vary.

\subsection{A finite tuple that exact algebras cannot realize}

The graph measurements are fixed before their entries are chosen.
To choose those entries under nuclearity, we use a finite tuple whose
linear matrix norms cannot all occur in an exact algebra, even though
each fixed such norm is a limit of norms of finite matrix tuples.
Nonexact examples originate in work of Wassermann
\cite{Wassermann1976}.  Pisier develops the operator-space formulation
and finite compression mechanism in
\cite[Theorems~1 and~7, Lemma~2 and Equation~(7)]{Pisier1995}.
We use the existence of a nonexact algebra from this theory and give
the tensor norm and compression arguments needed here.  They retain
the quantifiers across all matrix levels required for the graph entries.

Recall that a \(C^*\)-algebra \(D\) is \emph{exact} if tensoring every
short exact sequence with \(D\) in the minimal tensor norm preserves
exactness.  Thus exactness identifies the quotient norm on
\((D\otimes_{\min}J)/(D\otimes_{\min}I)\) with the spatial norm on
\(D\otimes_{\min}(J/I)\), for every closed ideal \(I\triangleleft J\).
The next proof starts from a strict inequality between these norms.

\begin{lemma}[A finite tuple and finite matrix approximations]
\label{lem:forbidden-tuple}
There are \(\ell\ge1\), a \(C^*\)-algebra \(W\), contractions
\(z_1,\ldots,z_\ell\in W\), and contractions
\(z_j^{(r)}\in M_{N_r}(\C)\), where \(N_r\ge1\), \(r\ge1\), and
\(1\le j\le\ell\), with the following properties.
\begin{enumerate}[label=\textup{(\roman*)}]
\item No tuple \(b_1,\ldots,b_\ell\) in any exact \(C^*\)-algebra
satisfies
\begin{equation}\label{eq:forbidden-tuple}
 \left\|\sum_{j=1}^{\ell} A_j\otimes b_j\right\|
 =\left\|\sum_{j=1}^{\ell} A_j\otimes z_j\right\|
 \quad\text{for every }m\ge1\text{ and every }A_j\in M_m(\C).
\end{equation}
\item For every fixed \(m\ge1\) and \(A_j\in M_m(\C)\),
\begin{equation}\label{eq:tuple-approximations}
 \lim_{r\to\infty}
 \left\|\sum_{j=1}^{\ell} A_j\otimes z_j^{(r)}\right\|
 =\left\|\sum_{j=1}^{\ell} A_j\otimes z_j\right\|.
\end{equation}
\end{enumerate}
\end{lemma}

\begin{proof}
Take a nonexact \(C^*\)-algebra \(W\), whose existence is recorded
in \cite[Remark~(i) following Theorem~1, pp.~163--164]{Pisier1995}.
There are
\(I\triangleleft J\) such that the canonical surjection
\[
 (W\otimes_{\min}J)/(W\otimes_{\min}I)
 \longrightarrow W\otimes_{\min}(J/I)
\]
is not injective.  The map is surjective because its algebraic image
is dense and the range of a \(C^*\)-homomorphism is closed.  Choose a
representative \(T_0\in W\otimes_{\min}J\) of a nonzero class in its
kernel.  Its norm modulo \(W\otimes_{\min}I\) is positive, whereas
its image in \(W\otimes_{\min}(J/I)\) is zero.  A sufficiently close
finite algebraic approximation to \(T_0\) retains a strict gap between
these two norms.  After normalizing its nonzero first factors and
absorbing their norms into the second factors, we obtain contractions
\(z_j\in W\) and elements \(h_j\in J\) for which
\begin{equation}\label{eq:tuple-gap}
 \left\|\sum_j z_j\otimes h_j+W\otimes_{\min}I\right\|
 >\left\|\sum_j z_j\otimes(h_j+I)\right\|_{\min}.
\end{equation}
For a tuple \(\mathbf v=(v_j)_{j=1}^{\ell}\) in a \(C^*\)-algebra
\(D\), write
\[
 \mu_{\mathbf v,m}(A_1,\ldots,A_\ell)
 =\left\|\sum_{j=1}^{\ell}A_j\otimes v_j\right\|_{M_m(D)}.
\]
We show that these functions determine both norms in
\eqref{eq:tuple-gap}, with the coefficient algebra and its ideal
held fixed.

For arbitrary coefficients \(d_j\in J\), represent the fixed
nonzero algebra \(J\) faithfully on a Hilbert space \(\mathcal H\).
Then
\begin{equation}\label{eq:tuple-compression-formula}
 \left\|\sum_j v_j\otimes d_j\right\|_{\min}
 =\sup_{m,V}
 \mu_{\mathbf v,m}(V^*d_1V,\ldots,V^*d_\ell V),
\end{equation}
where \(V:\C^m\to\mathcal H\) ranges over finite-dimensional
isometries.  Indeed, also represent \(D\) faithfully.  The operators
on the right are the two-sided compressions in the second Hilbert
factor, so their norms are at most the left-hand norm.  Finite-rank
projections, directed by inclusion, converge strongly to the
identity.  Their two-sided compressions converge strongly to the
tensor sum; evaluation at a fixed vector whose image nearly attains
its norm gives the reverse inequality.  This applies to arbitrary,
including nonselfadjoint, tensor sums.

For the quotient norm, choose a positive contractive approximate
identity \((f_\lambda)\) of the fixed ideal \(I\).  Put
\(J_D=D\otimes_{\min}I\).  For \(T\in D\otimes_{\min}J\), define
\[
 R_\lambda(T)
 =(1\otimes(1-f_\lambda))T(1\otimes(1-f_\lambda)),
\]
using the unitized multiplier operations.  Elementary tensors and
density give
\[
 T-R_\lambda(T)\in J_D,\qquad
 \|R_\lambda(K)\|\longrightarrow0\quad\text{for each fixed }K\in J_D.
\]
The cutdowns are contractions.  Hence, for every fixed \(K\in J_D\),
\[
 \dist(T,J_D)\le\|R_\lambda(T)\|
 \le\|T-K\|+\|R_\lambda(K)\|.
\]
Taking lower and upper limits and then the infimum over \(K\) yields
the net limit
\begin{equation}\label{eq:tuple-quotient-formula}
 \left\|\sum_j v_j\otimes d_j+J_D\right\|
 =\lim_\lambda
 \left\|\sum_j v_j\otimes
 (1-f_\lambda)d_j(1-f_\lambda)\right\|_{\min}.
\end{equation}
The cutdown coefficients lie in \(J\), and no quasicentrality is
needed.  This argument uses convergence on each fixed element of
\(J_D\).  A related fixed-element cutdown mechanism appears in
\cite[proof of Lemma~3, p.~165]{Pisier1995}; the calculation above
establishes the stated formula for general \(I\triangleleft J\) and
its approximate-identity net.  Equations~\eqref{eq:tuple-compression-formula} and
\eqref{eq:tuple-quotient-formula} show that the displayed quotient
norm depends only on the functions \(\mu_{\mathbf v,m}\).  Applying
the same compression argument in \(J/I\) shows the same for
\(\|\sum_j v_j\otimes(d_j+I)\|_{\min}\).

If a tuple in an exact algebra had the same functions
\(\mu_{\mathbf z,m}\) at every matrix level, the two formulas above
would transfer this strict gap to that algebra.  Exactness makes its
canonical quotient map isometric, a contradiction.  This proves
(i).  The conclusion excludes simultaneous equality for all finite
matrix tests; it does not assert that a fixed finite list of tests
suffices.

For (ii), set \(W_0=C^*(z_1,\ldots,z_\ell)\).  This algebra is
separable, and its inclusion in \(W\) preserves matrix norms.  It
has a faithful representation on a separable Hilbert space.  Indeed,
choose a dense sequence in the positive unit sphere and, for each
member, a state attaining its norm.  Their GNS representations have
faithful direct sum: for every nonzero positive element, approximation
by that sequence shows that some chosen state is positive on it.
Choose increasing nonzero
finite-rank projections \(P_r\to1\) strongly in that representation,
and set
\[
 N_r=\dim P_r\mathcal H,\qquad
 z_j^{(r)}=P_rz_jP_r\big|_{P_r\mathcal H}.
\]
These are contractions.  For each fixed \(m\) and \((A_j)\), the
corresponding two-sided compressions of \(\sum_j A_j\otimes z_j\)
converge strongly to that operator and have norms at most its norm.
Evaluation at a fixed near-norming vector gives the matching lower
limit, proving \eqref{eq:tuple-approximations}.  The same sequence
works for every fixed test.  The assertion is pointwise in the test,
and uses no multiplicativity of the compressions.
\end{proof}

\subsection{Choosing graph entries under nuclearity}

We use the finite tuples from Lemma~\ref{lem:forbidden-tuple} to
choose entries in a spectral hereditary subalgebra.  The central
point is that the transport lemma provides one rectangular map for
the whole tuple.  Its conjugation preserves the norm of every joint linear
matrix expression, not merely the norms of the individual entries.

\begin{proposition}[Nuclear avoidance of the fixed measurements]
\label{prop:nuclear-avoidance}
Let \(C\) be nuclear and satisfy \(\propF\) for a fixed \(n\ge1\),
and let \(\ell\) be as in Lemma~\ref{lem:forbidden-tuple}.  For
\(X\in C_+\) and fixed elements
\(B_{j,pq}\in C\), where \(1\le j\le\ell\) and \(1\le p,q\le n\),
there are contractions
\[
 S_{j,pq}\in\Her((X-\alpha_0)_+)
\]
such that no representation \(\pi\) satisfies simultaneously
\begin{equation}\label{eq:nuclear-forbidden}
 \|\pi(X)\|>s,\qquad
 \pi(B_{j,pq})=\pi(S_{j,pq})
 \quad(1\le j\le\ell,\ 1\le p,q\le n).
\end{equation}
Here \(\alpha_0=1/32\) and \(s=1/4\), as in
Proposition~\ref{prop:graph}.
\end{proposition}

\begin{proof}
Put \(g=(X-\alpha_0)_+\), choose \(0<\nu<s-\alpha_0\), and put
\(h_g=(g-\nu)_+\).  Choose a ramp \(t'=t'(X)\) that is zero on
\([0,\alpha_0+\nu]\) and one on \([s,\infty)\).  Then
\(t'\in\Her(h_g)\), and
\begin{equation}\label{eq:nuclear-tuple-signal}
 \|\pi(X)\|>s\quad\Longrightarrow\quad\|\pi(t')\|=1.
\end{equation}
The element \(v=g^{\oplus n}\in M_n(C)\) is properly infinite by
\(\propF\), applied to the scalar positive \(g\).  For the dimensions
\(N_r\) of Lemma~\ref{lem:forbidden-tuple},
\[
 d_r=g^{\oplus N_r}\precsim v^{\oplus N_r}\precsim v.
\]
The first comparison uses \(g\) as a summand of \(v\); the second
uses proper infiniteness of \(v\).  Lemma~\ref{lem:transport}, with
cutoff \(\nu\), gives one rectangular carrier
\[
 w_r\in M_{n,N_r}(C^{**}),\qquad
 w_r^*w_r=1_{N_r}\otimes e_{h_g},\qquad
 w_rw_r^*\le1_n\otimes e_g.
\]
Its conjugation sends \(M_{N_r}(\Her(h_g))\) into
\(M_n(\Her(g))=\Her_{M_n(C)}(v)\).  Use this same carrier for every
\(j\), defining
\begin{equation}\label{eq:nuclear-tuple-choices}
 S_j^{(r)}=w_r(z_j^{(r)}\otimes t')w_r^*
 \in M_n(\Her(g)),\qquad
 S_{j,pq}^{(r)}=(S_j^{(r)})_{pq}.
\end{equation}
Each matrix has norm at most one, so its entries are contractions in
the required hereditary subalgebra.

Suppose that every \(r\) admits a representation \(\pi_r\) satisfying
\eqref{eq:nuclear-forbidden} for these entries.  Restricting to a
reducing cyclic subrepresentation that still detects
\(\|\pi_r(X)\|>s\), we may make these choices from a set of
representations.  Write
\(B_j=[B_{j,pq}]\in M_n(C)\); these elements are fixed independently
of \(r\).  At every matrix level \(m\), the normal amplification of
the single \(w_r\) is isometric on the supported source algebra, by
Lemma~\ref{lem:transport}.  The represented entry equalities and
\eqref{eq:nuclear-tuple-signal} therefore give, for every
\(A_j\in M_m(\C)\),
\begin{align}
 \left\|\sum_j A_j\otimes\pi_r^{(n)}(B_j)\right\|
 &=\left\|\left(\sum_j A_j\otimes z_j^{(r)}\right)
       \otimes\pi_r(t')\right\|\notag\\
 &=\left\|\sum_j A_j\otimes z_j^{(r)}\right\|.
 \label{eq:nuclear-exact-matrix-norms}
\end{align}
This equality is exact at each \(r\).  Using a different carrier for
each \(j\) would not give the joint equality.

Define the \(*\)-homomorphism
\[
 \Psi:M_n(C)\longrightarrow
 \mathcal Q=
 \frac{\prod_r B(\mathcal H_r^{\oplus n})}
      {\bigoplus_r B(\mathcal H_r^{\oplus n})},
 \qquad
 \Psi(b)=[(\pi_r^{(n)}(b))_r],
\]
where the numerator is the algebra of bounded sequences and the
denominator is its ideal of norm-null sequences.  Contractivity of
the representations makes each displayed sequence bounded.  The
image \(E=\Psi(M_n(C))\) is closed and isomorphic to
\(M_n(C)/\ker\Psi\).  Since \(M_n(C)\cong M_n\otimes C\), nuclearity
of matrix algebras and permanence under tensor products make
\(M_n(C)\) nuclear \cite[II.9.4.2 and IV.3.1.1]{Blackadar}.
Its quotient \(E\) is nuclear by
\cite[Corollary~4]{ChoiEffros1977}, and hence exact by
\cite[II.9.6.6]{Blackadar}.
These permanence facts require neither units nor separability.
It is this quotient of a nuclear algebra that must be exact; the
argument requires no exactness of the ambient \(\mathcal Q\).

The norm of a class modulo norm-null sequences is the \(\limsup\) of
its coordinate norms.  Subtracting a norm-null sequence gives the
lower bound, and subtracting finitely many initial entries gives the
upper bound.  The same formula holds at each fixed matrix level under
the canonical identifications of finite matrices over a product and
its norm-null ideal.  Since \(E\subset\mathcal Q\) is completely
isometric, \eqref{eq:nuclear-exact-matrix-norms} and
\eqref{eq:tuple-approximations} give
\[
 \left\|\sum_j A_j\otimes\Psi(B_j)\right\|_{M_m(E)}
 =\limsup_r\left\|\sum_j A_j\otimes\pi_r^{(n)}(B_j)\right\|
 =\left\|\sum_j A_j\otimes z_j\right\|
\]
for every \(m\) and every \((A_j)\).  Only the fixed elements
\(B_j\) have been sent into \(E\); the varying choices in
\eqref{eq:nuclear-tuple-choices} were used in the coordinatewise
equalities.  The displayed identity contradicts
Lemma~\ref{lem:forbidden-tuple}(i).  Thus some \(r\) admits no such
representation, and its entries in
\eqref{eq:nuclear-tuple-choices} prove the proposition.
\end{proof}

\begin{proposition}[Shrinking two corners under nuclearity]
\label{prop:shrinking-nuclear}
Let \(C\) be nuclear and satisfy \(\propF\) for one fixed \(n\ge1\).
Let \(a,y\in C_+\) be contractions, with \(a\) rich, and let
\(c\in\widetilde C\) be a contraction.  There are positive
contractions \(p_1\in\Her(a)\) and \(q_1\in\Her(y)\) such that
\(p_1\) is rich,
\begin{equation}\label{eq:nuclear-shrinking-signals}
 \begin{split}
 \|\pi(a)\|>1/2&\quad\Longrightarrow\quad\|\pi(p_1)\|=1,\\
 \|\pi(y)\|>1/2&\quad\Longrightarrow\quad\|\pi(q_1)\|=1
 \end{split}
 \qquad\text{for every representation }\pi,
\end{equation}
and \(\|e_{p_1}ce_{q_1}\|_{C^{**}}<1\).
\end{proposition}

\begin{proof}
Choose \(\ell\) from Lemma~\ref{lem:forbidden-tuple}, set
\(k=\ell n^2\), and choose graph coefficients
\(\gamma_g,\sigma_g>0\) with \(\gamma_g^2+k\sigma_g^2<1\).
With these choices fixed, apply Proposition~\ref{prop:graph} to obtain
its measurements \(X,B_i\).  Group the \(B_i\) into \(\ell\) matrices
\(B_j=[B_{j,pq}]\in M_n(C)\), and choose the corresponding
\(S_{j,pq}\) by Proposition~\ref{prop:nuclear-avoidance}.
The graph gives a rich positive contraction \(p\in\Her(a)\) with
\(\|\pi(p)\|=1\) whenever \(\|\pi(a)\|>1/2\).  If
\(d=y(1-c^*p^2c)y\) vanishes in a representation where
\(\|\pi(y)\|>s\), the graph would give all of
\eqref{eq:nuclear-forbidden}, contrary to the choice of the entries.
Thus
\begin{equation}\label{eq:nonvanishing-N}
 \pi\bigl(y(1-c^*p^2c)y\bigr)=0
 \quad\Longrightarrow\quad\|\pi(y)\|\le s=1/4
\end{equation}
for every representation.  Lemma~\ref{lem:conditional-shrinking}
applies to this \(p\), its norm signal, and the displayed implication,
and gives the asserted \(p_1,q_1\).  The one fixed \(n\) controlled
both the number of entries and the algebra \(M_n(C)\) in the quotient
argument; only the finite dimensions \(N_r\) varied.
\end{proof}

\subsection{Local rich elements and their merger}

The shrinking proposition will let us join rich elements once they
are available.  For \(u\in C_+\), we first produce such an element inside \(\Her(u)\) at
each irreducible representation that detects \(u\).  This local
construction needs only the fixed-amplification hypothesis
\(\propF\); nuclearity will enter when the elements are joined.

\begin{lemma}[A rich element detected by one representation]
\label{lem:local-rich}
Let \(C\) be a \(C^*\)-algebra, and suppose that one integer
\(n\ge1\) makes \(h^{\oplus n}\) properly infinite for every
\(h\in C_+\).  If \(u\in C_+\) and an irreducible representation
\(\pi:C\to B(\mathcal H)\) satisfies \(\pi(u)\ne0\), then there is
a rich positive contraction \(p\in\Her(u)\) with \(\|\pi(p)\|=1\).
\end{lemma}

\begin{proof}
We will find \(n\) orthogonal targets in \(\Her(u)\), each detected
by \(\pi\), and one positive element comparing into every target.
Transporting \(n\) copies of its cutoff will allow \(\propF\) to
produce richness inside \(\Her(u)\).

Put \(\mathcal H_u=\overline{\pi(u)\mathcal H}\).  This space is
infinite dimensional.  Otherwise \(\pi(u)\) would have finite rank
\(r>0\).  Applying the stabilized representation to proper
infiniteness of \(u^{\oplus n}\) would approximate
\(\pi(u)^{\oplus2n}\), of rank \(2nr\), by positive operators of
rank at most \(nr\).  On the \(2nr\)-dimensional range of the target,
each such approximant has a unit vector in its kernel.  Its distance
from the target is therefore at least the least positive eigenvalue
of \(\pi(u)\), contradicting norm approximation.

The action of \(D=\Her(u)\) on \(\mathcal H_u\) is irreducible.
Indeed, \(\pi(u)\) has no kernel on its support, so
\(\pi(u)\xi\ne0\) for \(0\ne\xi\in\mathcal H_u\).
Irreducibility of \(\pi\) gives
\(\overline{\pi(C)\pi(u)\xi}=\mathcal H\).  Applying \(\pi(u)\)
and taking closure yields
\[
 \overline{\pi(uCu)\xi}=\mathcal H_u.
\]
Thus every nonzero vector is cyclic for the restricted hereditary
algebra, which proves irreducibility.

Apply the selfadjoint form of Kadison transitivity
\cite[Theorem~1]{Kadison1957} to the irreducible unitization of this
action.  Choose \(n+1\) orthonormal vectors and prescribe \(n+1\)
distinct real eigenvalues.  A selfadjoint interpolant has the form
\(h_0+\lambda1\), with \(h_0\in D_{\rm sa}\).  The corresponding
eigenvalues of \(\pi(h_0)\) remain distinct, so at most one is zero.
Choose \(n\) nonzero ones.  Nonnegative continuous functions
supported in pairwise disjoint neighborhoods of these values,
avoiding zero and positive at the chosen values, give orthogonal
elements
\[
 b_1,\ldots,b_n\in D_+,\qquad \pi(b_i)\ne0.
\]
The functions vanish at zero, so these elements belong to \(D\)
also when \(D\) is nonunital.

There is an element \(b'\in D_+\), still detected by \(\pi\), with
\(b'\precsim b_i\) for every \(i\).  Start with \(e'=b_1\).
Given the current \(e'\) and the next \(b_i\), choose a vector on
which \(\pi(e')\) is nonzero.  Cyclicity for \(D\) and nonzero
\(\pi(b_i)\) give \(h\in D\) such that \(\pi(z)\ne0\), where
\(z=b_ihe'\).  Then
\[
 z^*z\in\Her(e'),\qquad zz^*\in\Her(b_i),\qquad
 z^*z\CuEq zz^*.
\]
Replacing \(e'\) by \(z^*z\) preserves nonzero image and every
previous comparison, by transitivity.  At the end we obtain \(b'\)
and
\begin{equation}\label{eq:local-rich-comparison}
 (b')^{\oplus n}\precsim b_1\oplus\cdots\oplus b_n
 \CuEq\sum_{i=1}^{n}b_i\precsim u.
\end{equation}
The middle equivalence uses orthogonality, and the last comparison
uses \(\sum_i b_i\in\Her(u)\).

Choose \(\tau>0\) with \(2\tau<\|\pi(b')\|\).  Transport the
\(n\) source copies in \eqref{eq:local-rich-comparison}, after cutoff
\(\tau\), into \(\Her(u)\).  Let \(t\) be a ramp of \(b'\), zero
on \([0,\tau]\) and one on \([2\tau,\infty)\).  It belongs to the
common cutoff hereditary algebra and has represented norm one.
The sum \(p\) of its \(n\) orthogonal transported copies is a
positive contraction in \(\Her(u)\) with \(\|\pi(p)\|=1\).
By \eqref{eq:orthogonal-copy-cutoffs}, for every \(\varepsilon\ge0\),
\[
 (p-\varepsilon)_+\CuEq((t-\varepsilon)_+)^{\oplus n}.
\]
The right side is properly infinite by the hypothesis applied to
the scalar positive \((t-\varepsilon)_+\).  Thus \(p\) is rich.
\end{proof}

We now join two rich elements while retaining the union of their
norm signals.  At each step Proposition~\ref{prop:shrinking-nuclear}
is used with \(c=1\) in the unitization.  Its support bound is one
global norm below one for that step; the numerical bound may change
at the next step.

\begin{lemma}[Merging two rich elements]\label{lem:merger}
Let \(C\) be nuclear and satisfy \(\propF\) for a fixed \(n\ge1\).
For \(u\in C_+\) and rich positive contractions \(a,b\in\Her(u)\),
there is a rich positive contraction \(c_0\in\Her(u)\) such that
every representation \(\pi\) satisfies
\begin{equation}\label{eq:merger-signal}
 \bigl(\|\pi(a)\|>1/2\ \text{or}\ \|\pi(b)\|>1/2\bigr)
 \quad\Longrightarrow\quad\|\pi(c_0)\|=1.
\end{equation}
\end{lemma}

\begin{proof}
All comparisons remain in \(C\otimes\Kc\), and all uses of shrinking
are in \(C\).  Since \(b\) is properly infinite,
\(b^{\oplus n}\precsim b\).  Apply Lemma~\ref{lem:transport} with
cutoff \(1/8\) to obtain orthogonal copy maps
\[
 \theta_i(h)=w_i h w_i^*:
 \Her((b-1/8)_+)\longrightarrow\Her(b)\subset\Her(u)
 \quad(1\le i\le n).
\]
Let \(y_0\) be a ramp of \(b\), zero on \([0,1/4]\) and one on
\([1/2,1]\), and set \(a_0=a\).  For \(1\le j\le n\) we construct
positive contractions \(a_j,y_j\) such that \(a_j\) is rich and
belongs to \(\Her(u)\), \(y_j\) belongs to the common domain of
the copy maps, and the following invariants hold:
\begin{equation}\label{eq:merger-induction}
 \begin{gathered}
 a_j\theta_i(y_j)=0\qquad(1\le i\le j),\\
 \|\pi(a)\|>1/2\ \Longrightarrow\ \|\pi(a_j)\|=1,\\
 \|\pi(b)\|>1/2\ \Longrightarrow\ \|\pi(y_j)\|=1
 \quad\text{for every representation }\pi.
 \end{gathered}
\end{equation}

Once the induction reaches \(j=n\), we will set
\(c_0=a_n+\sum_{i=1}^n\theta_i(y_n)\).  The induction will make the
summands orthogonal, and the final use of \(\propF\) will prove the
sum of the \(n\) copies rich without requiring \(y_n\) itself to be
rich.

Suppose \(a_{j-1},y_{j-1}\) have been constructed.  At \(j=1\),
use the initial pair \(a_0,y_0\); at later stages use the preceding
invariants.  Apply
Proposition~\ref{prop:shrinking-nuclear} to
\(a_{j-1},\theta_j(y_{j-1})\), and \(1\in\widetilde C\), and call
its outputs \(p_j,q_j\).  Thus \(p_j\) is rich and belongs to
\(\Her(a_{j-1})\), while
\[
 \kappa_j=\|e_{p_j}e_{q_j}\|<1.
\]
The hereditary-image equality
\eqref{eq:transport-hereditary-image} supplies a positive contraction
\(z_j\in\Her(y_{j-1})\) such that
\begin{equation}\label{eq:merger-pullback}
 q_j=\theta_j(z_j).
\end{equation}
This pullback is an algebra element.  Its represented norms equal
those of \(q_j\) by \eqref{eq:transport-represented-norm}.

For \(0\le t\le1\), put \(f(t)=(2t-1)_+\) and
\(g(t)=\min\{1,2t\}\).  We next move \(p_j\) away from a cutoff
of \(q_j\).  The support
bound ensures that multiplication by \(1-g(q_j)\) is bounded below
on the support of \(p_j\).  The polar part of
\((1-g(q_j))p_j\) will therefore transport \(p_j\), preserving its
cutoff Cuntz classes and represented norms.
Define
\[
 y_j=f(z_j),\qquad H_j=1-g(q_j),\qquad T_j=H_jp_j.
\]
Here \(T_j=p_j-g(q_j)p_j\in\Her(u)\).  In the universal
representation, \(g(q_j)\) is supported under \(e_{q_j}\).
For \(\xi\in e_{p_j}\mathcal H\), this gives
\[
 \|H_j\xi\|^2
 \ge\|(1-e_{q_j})H_j\xi\|^2
 =\|(1-e_{q_j})\xi\|^2
 \ge(1-\kappa_j^2)\|\xi\|^2.
\]
Since \(H_j\) is a contraction, we obtain the order bounds
\begin{equation}\label{eq:merger-polar-order}
 (1-\kappa_j^2)p_j^2\le T_j^*T_j\le p_j^2.
\end{equation}
Let \(T_j=l_j|T_j|\) be its polar decomposition in the bidual.
The order bounds give \(\Her(|T_j|)=\Her(p_j)\) and
\(l_j^*l_j=e_{p_j}\).  The polar membership argument in the proof
of Lemma~\ref{lem:transport}, applied now to \(T_j\), gives
\[
 a_j=l_jp_jl_j^*\in\Her(T_jT_j^*)\subset\Her(u).
\]
It is a positive contraction.  The same argument makes each
transported cutoff Cuntz equivalent to the corresponding cutoff of
\(p_j\), so \(a_j\) is rich.  The normal extensions of the
representations preserve \(l_j^*l_j=e_{p_j}\); conjugating back by
\(l_j^*\) shows that \(a_j\) and \(p_j\) have equal norm in every
representation.

We check that this change preserves every earlier orthogonality
relation.  Since \(gf=f\), functional calculus and
\eqref{eq:merger-pullback} give
\[
 \theta_j(y_j)H_j=f(q_j)(1-g(q_j))=0,
\]
and hence \(\theta_j(y_j)T_j=0\).  For \(i<j\), the previous
orthogonality and \(p_j\in\Her(a_{j-1})\) give
\(\theta_i(y_{j-1})p_j=0\), while orthogonality of the copy ranges
gives \(\theta_i(y_{j-1})g(q_j)=0\).  Thus
\(\theta_i(y_{j-1})T_j=0\).  Each of these positive elements
annihilates \(\overline{\ran(T_j)}\), whose projection is
\(l_jl_j^*\), and so annihilates \(a_j\).  Because
\(y_j\in\Her(y_{j-1})\), the same is true of \(\theta_i(y_j)\)
for \(i<j\).  This proves the orthogonality in
\eqref{eq:merger-induction}.

The norm signals also persist simultaneously.  Wherever
\(\|\pi(a)\|>1/2\), the first shrinking step gives
\(\|\pi(p_1)\|=1\), hence \(\|\pi(a_1)\|=1\); that signal is
above \(1/2\) and survives every subsequent step.  Wherever
\(\|\pi(b)\|>1/2\), the initial ramp has \(\|\pi(y_0)\|=1\).
If \(\|\pi(y_{j-1})\|=1\), its copy under \(\theta_j\) has norm
one, shrinking gives \(\|\pi(q_j)\|=1\), and transport gives
\(\|\pi(z_j)\|=1\).  Since \(f(1)=1\), spectral mapping gives
\(\|\pi(y_j)\|=1\).

Finally set
\[
 c_0=a_n+\sum_{i=1}^{n}\theta_i(y_n).
\]
All summands are orthogonal positive contractions in \(\Her(u)\).
Their sum is a contraction, and in each representation its norm is
the maximum of the summand norms.  The preceding signals prove
\eqref{eq:merger-signal}.  For every \(\tau\ge0\), orthogonal
functional calculus and \eqref{eq:orthogonal-copy-cutoffs} give
\[
 (c_0-\tau)_+
 \CuEq(a_n-\tau)_+\oplus((y_n-\tau)_+)^{\oplus n}.
\]
The first summand is properly infinite because \(a_n\) is rich.
The second is properly infinite by \(\propF\), applied to the
scalar positive \((y_n-\tau)_+\).  Their block sum is properly
infinite, so \(c_0\) is rich.  The construction includes zero
inputs through zero copy maps and zero cutoffs.
\end{proof}

\subsection{Finite ideal assembly and absorption}

Lemma~\ref{lem:local-rich} supplies a rich element wherever \(u\) is
detected, and Lemma~\ref{lem:merger} joins any finite family while
retaining all of its norm signals.  These two facts put
\((u-\varepsilon)_+\), for every \(\varepsilon>0\), in the ideal of
a single rich element.  Ideal comparison then supplies the doubled
cutoff comparison needed for Theorem~\ref{thm:nuclear-to-individual}.

\begin{proof}[Proof of Theorem~\ref{thm:nuclear-to-individual}]
Fix \(u\in C_+\), and let \(\mathcal R_u\) be the set of rich positive
contractions in \(\Her(u)\).  Let \(J\) be the closed ideal generated
by all \((p-1/2)_+\) with \(p\in\mathcal R_u\).  We claim that
\(u\in J\).  Otherwise an irreducible representation of \(C/J\)
would detect \(u+J\).  Its composition with the quotient map is an
irreducible representation of \(C\), so Lemma~\ref{lem:local-rich}
gives \(p\in\mathcal R_u\) with represented norm one.  Then
\((p-1/2)_+\) has nonzero image, contrary to its membership in \(J\).

The ideals generated by finite subfamilies of these cutoffs form an
increasing directed family with dense union in \(J\).  For
\(\varepsilon>0\), Lemma~\ref{lem:ideals-cutoffs}(i) therefore puts
\((u-\varepsilon)_+\) in the ideal generated by
\((p_i-1/2)_+\) for finitely many \(p_1,\ldots,p_r\in\mathcal R_u\).
Starting from the rich zero element, merge them successively by
Lemma~\ref{lem:merger}.  The result is a rich \(p\in\Her(u)\) for
which every representation satisfies
\[
 \|\pi(p_i)\|>1/2\text{ for some }i
 \quad\Longrightarrow\quad\|\pi(p)\|=1.
\]
A norm-one signal survives the next merger because it is above
\(1/2\).  The primitive ideal test in
Lemma~\ref{lem:ideals-cutoffs}(ii) now gives
\((p_i-1/2)_+\in\Id(p)\) for every \(i\).  Consequently
\begin{equation}\label{eq:nuclear-final-ideal}
 (u-\varepsilon)_+\in\Id(p).
\end{equation}
If the cutoff is zero, the zero element already suffices.

Since \(p\) is properly infinite, ideal comparison with this target in
the stabilization, followed by hereditary comparison, yields
\begin{equation}\label{eq:nuclear-final-comparison}
 ((u-\varepsilon)_+)^{\oplus2}\precsim p\precsim u.
\end{equation}
The target \(u\) is fixed as \(\varepsilon\downarrow0\).  Explicitly,
choose \(\varepsilon_j\downarrow0\) and, for each \(j\), a comparison
coefficient in \eqref{eq:nuclear-final-comparison} with error less
than \(\varepsilon_j\) to the doubled cutoff.  That cutoff differs
from \(u\oplus u\) by at most \(\varepsilon_j\).  The resulting
coefficients therefore compare \(u\oplus u\) to \(u\) in norm.
Thus \(u\) is properly infinite.  The fixed \(n\) was applied only
to scalar positives and their scalar cutoffs throughout the proof.
\end{proof}

\begin{remark}[Extension of the tuple proof to exact algebras]
\label{rem:tuple-exact}
The amplification-removal argument uses nuclearity only to obtain
exactness of \(E=\Psi(M_n(C))\) in
Proposition~\ref{prop:nuclear-avoidance}.  Quotients of exact
\(C^*\)-algebras are exact, including in the nonseparable case
\cite[IV.3.4.20]{Blackadar}.  Finite matrix amplification also
preserves exactness, since its tensor products and ideal kernels
are the corresponding finite matrices over those for \(C\).
This deeper quotient-permanence theorem therefore makes the same
avoidance, shrinking, merger, and ideal assembly valid when \(C\)
is merely exact.  It gives another proof of the amplification-removal
conclusion of Theorem~\ref{thm:exact-infiniteness}.  The perturbation
and localization proof below obtains that conclusion without invoking
quotient permanence of exactness.
\end{remark}

\begin{proof}[Proof of Corollary~\ref{cor:absorption}]
Set \(C=A\), the separable nuclear algebra in the corollary, with
its stated fixed-amplification hypothesis.
Theorem~\ref{thm:nuclear-to-individual} gives \(\propP\), and
Theorem~\ref{thm:individual-to-strong} gives strong pure infiniteness.
The general nonunital absorption implication is recorded in the
paragraph following \cite[Definition~3.1]{KS}.  We spell out the
central-sequence interface that includes the nonunital case.

Let \(C_\omega\) be the norm ultrapower of \(C\) along a free
ultrafilter on \(\N\), with \(C\) embedded diagonally, and put
\[
 F(C)=\frac{C_\omega\cap C'}{\operatorname{Ann}(C,C_\omega)},
 \qquad
 \operatorname{Ann}(C,C_\omega)
 =\{x\in C_\omega:xC=Cx=0\}.
\]
Suppose first that \(C\ne0\).  Separability makes \(C\)
\(\sigma\)-unital, so \(F(C)\) is unital by
\cite[Corollary~1.10(1)]{Kirchberg2006}.  The central-sequence embedding
statement immediately before \cite[Proposition~4.4]{Kirchberg2006}
gives a unital embedding \(\Oinf\hookrightarrow F(C)\): its ambient
original algebra is the strongly purely infinite \(C\), and the
diagonal copy of \(C\) is the required separable nuclear subalgebra
of \(C_\omega\).  The algebra \(\Oinf\) is unital and separable, and
Remark~4.6 supplies self-absorption and connectedness of its unitary
group.  The ``in particular'' clause of Proposition~4.4(4) therefore
gives stable \(\Oinf\)-absorption.  Since \(U(\Oinf)=U_0(\Oinf)\), where \(U_0\)
is the identity component, the commutator subgroup of \(U(\Oinf)\)
lies in \(U_0(\Oinf)\); part (5) therefore removes stabilization.  Hence
\(C\cong C\otimes_{\min}\Oinf\).  The Proposition permits \(C\) to
be nonunital; the unital embedding is into the annihilator quotient
\(F(C)\).  If \(C=0\), the asserted isomorphism is immediate.
\end{proof}
